\begin{table}[!h]
\centering
\caption{Table B2. Percentage-Percentage Method Comparison (Lundgren et al. 2018)}
\centering
\resizebox{\ifdim\width>\linewidth\linewidth\else\width\fi}{!}{
\fontsize{10}{12}\selectfont
\begin{threeparttable}
\begin{tabular}[t]{llllll}
\toprule
Metric & Pct-Count (restricted) & Pct-Pct Aggregate & Pct-Pct Core & Pct-Pct Peripheral & Lundgren et al. range\\
\midrule
N & 130 & 624 & 117 & 507 & \\
Mean & -0.386319 & -0.000000 & 0.015670 & -0.003616 & \\
SD & 1.4008 & 0.1848 & 0.3505 & 0.1176 & \\
Raw Kurtosis (benchmark=3) & 37.9179 & 17.8077 & 5.2949 & 39.6513 & \\
L-kurtosis (benchmark=0.1226) & 0.3721 & 0.6697 & 0.3131 & 0.7708 & EU=0.26, OAS=0.26, UN=0.28, AU=0.30, OIC=0.31\\
\addlinespace
Shapiro-Wilk W & 0.465710 & 0.582240 & 0.917113 & 0.427920 & \\
Shapiro-Wilk p & 9.6292e-20 & 4.1727e-36 & 2.1339e-06 & 3.7313e-37 & \\
Proportion zeros & 0.1462 & 0.7083 & 0.2564 & 0.8126 & \\
\bottomrule
\end{tabular}
\begin{tablenotes}[para]
\item \textit{Note: } 
\item Pct-pct method uses first differences of proportional shares. Lundgren et al. (2018) L-kurtosis values for comparison IOs.
\end{tablenotes}
\end{threeparttable}}
\end{table}
